A physically matched design comparison changes both the interpretation and selection of corrugated profiles. In the archived nonlinear campaign, actual medium volume per footprint reverses an asserted material saving and reduces the historically eligible C038 material--stress front from 31 members to S6 alone. Analytic knot-sided curvature rejects 16 of the 572 previously admitted geometries. These effects are consequences of resource definitions and geometric evaluation, rather than optimizer choice.

The new direct-FEM experiment fixes pitch, height and medium volume per initial area, then compares simple smooth references with NURBS. At the tested 64-path budget, three paired runs provide no consistent NSGA-II advantage over Sobol. The richer family expands the observed energy--stress trade-off, while the sinusoid remains a competitive reference.

Finite-scenario qualification changes the shortlist: four of ten candidates remain admissible after 27 perturbations, and enrichment removes a candidate that passed all axial tests. At 24 elements, a qualified energy-focused NURBS profile retains a 7.91\% sampled minimum-potential advantage over the sine with a 3.38\% sampled maximum-stress cost. Finer checks preserve the sign of this energy--stress trade-off in the inspected states, while its magnitude remains resolution-dependent. The nominal minimum-stress design and knee do not survive the geometric checks. Explicit resources, analytic geometric screening and enriched perturbation tests therefore provide the evidence needed to choose which profiles merit controlled physical prototyping.
